\chapter*{Lista de Termos e Abreviações}
\addcontentsline{toc}{chapter}{Lista de Termos e Abreviações}

\begin{description}
    \item[API (Application Programming Interface)] Interface de Programação de Aplicações. Um conjunto de regras e protocolos que permite a comunicação e a troca de dados entre diferentes sistemas de software.
    \item[Automação Residencial] Utilização de tecnologia para automatizar e controlar dispositivos e sistemas domésticos (iluminação, climatização, segurança), visando melhorar o conforto, a eficiência energética e a segurança.
    \item[BMS (Battery Management System)] Sistema de Gerenciamento de Bateria. Sistema eletrônico que gerencia uma bateria recarregável, protegendo-a contra operações inseguras, monitorando seu estado e otimizando seu desempenho.
    \item[CEP (Controle Estatístico de Processo)] Metodologia de qualidade que utiliza ferramentas estatísticas, como gráficos de controle, para monitorar e controlar a variabilidade de um processo, garantindo sua estabilidade e eficiência.
    \item[CUSUM (Cumulative Sum)] Soma Cumulativa. Técnica de análise sequencial e tipo de gráfico de controle altamente sensível para detectar pequenas, mas persistentes, mudanças na média de um processo.
    \item[DAG (Directed Acyclic Graph)] Grafo Acíclico Dirigido. Estrutura de dados que representa dependências entre tarefas em um fluxo de trabalho, onde cada nó é uma tarefa e uma aresta direcionada representa uma dependência.
    \item[Dagster] Ferramenta de orquestração de dados de código aberto projetada para construir, testar e monitorar pipelines de dados, com foco em ativos de dados e suas dependências.
    \item[Edge Computing (Computação de Borda)] Paradigma de computação distribuída que aproxima o processamento e o armazenamento de dados da fonte onde são gerados, a fim de melhorar o tempo de resposta e economizar largura de banda.
    \item[Estrutura Multicanal] Abordagem metodológica em que o comportamento de um sistema é modelado por meio de múltiplos canais de dados paralelos, cada um representando um período de tempo ou condição específica, permitindo uma análise estatística mais precisa.
    \item[ETL/ELT (Extract, Transform, Load / Extract, Load, Transform)] Processos de integração de dados para mover dados de múltiplas fontes para um repositório centralizado, como um data warehouse.
    \item[HMAC-SHA256] Código de Autenticação de Mensagem baseado em Hash usando SHA-256. Mecanismo criptográfico utilizado para verificar a integridade e a autenticidade de uma mensagem.
    \item[Interoperabilidade] Capacidade de diferentes sistemas, dispositivos ou aplicações de se conectarem e trocarem informações de maneira coordenada, sem restrições de acesso ou implementação.
    \item[IoT (Internet of Things)] Internet das Coisas. Rede de objetos físicos ("coisas") equipados com sensores, software e outras tecnologias para se conectar e trocar dados com outros dispositivos e sistemas pela internet.
    \item[JSON (JavaScript Object Notation)] Formato de intercâmbio de dados leve, de fácil leitura para humanos e de fácil interpretação para máquinas, amplamente utilizado em APIs web.
    \item[LGPD (Lei Geral de Proteção de Dados)] Legislação brasileira que regula a coleta, o uso, o processamento e o armazenamento de dados pessoais, com o objetivo de proteger os direitos e a privacidade dos indivíduos.
    \item[NAS (Network-Attached Storage)] Dispositivo de armazenamento de dados dedicado e conectado a uma rede, que permite o acesso centralizado a arquivos por múltiplos usuários e dispositivos heterogêneos.
    \item[Não invasiva] No contexto tecnológico, refere-se a métodos ou dispositivos que podem ser instalados e operados sem exigir modificações estruturais em um sistema ou ambiente existente.
    \item[OLAP (Online Analytical Processing)] Sistema projetado para realizar análises multidimensionais em alta velocidade sobre grandes volumes de dados, suportando a tomada de decisão complexa.
    \item[OLTP (Online Transactional Processing)] Sistema que gerencia aplicações orientadas a um grande volume de transações curtas e rápidas, como inserções, atualizações e exclusões de dados.
    \item[Parquet] Formato de armazenamento de dados em colunas, de código aberto, otimizado para consultas analíticas eficientes e compressão de dados em larga escala.
    \item[PCB (Printed Circuit Board)] Placa de Circuito Impresso. Estrutura que suporta mecanicamente e conecta eletricamente componentes eletrônicos usando trilhas condutivas.
    \item[Plug and Play] Padrão que permite que um sistema de computador reconheça e configure automaticamente um novo dispositivo de hardware, sem a necessidade de intervenção manual do usuário.
    \item[Streamlit] Framework de código aberto em Python que simplifica a criação e o compartilhamento de aplicações web interativas para projetos de ciência de dados e aprendizado de máquina.
    \item[UI (User Interface)] Interface do Usuário. Refere-se à parte visual e interativa de um sistema ou aplicação com a qual o usuário interage diretamente. Abrange o design dos elementos gráficos, como botões, ícones, menus e layout, focando na estética e na clareza da apresentação das informações e controles.
    \item[UX (User Experience)] Experiência do Usuário. Engloba a jornada completa e os sentimentos de um usuário ao interagir com um produto ou serviço. Vai além da interface visual, incluindo a usabilidade, a acessibilidade, o desempenho e a satisfação geral, com o objetivo de tornar a interação eficiente, eficaz e agradável.
\end{description}
